\begin{figure}[ht]
\centering
\begin{tikzpicture}[x=1cm, y=1cm, font=\small]

% Bi-objective design space: cost on x, mass on y. The achievable
% region is a blob; its lower-left boundary is the Pareto front.
% Front curve: y = 0.6 + 6/(x+0.2) on the window x in [1.0, 5.6],
% which gives y from ~5.6 down to ~1.6, kept inside the canvas.
\draw[->, thick] (-0.2, 0) -- (6.4, 0)
  node[right, font=\scriptsize] {objective 1: cost};
\draw[->, thick] (0, -0.2) -- (0, 4.6)
  node[above, font=\scriptsize] {objective 2: mass};

% Clip the shaded achievable region to a tidy box so it does not
% collide with the vertical axis.
\begin{scope}
  \clip (1.0, 0.0) rectangle (6.0, 4.3);
  \fill[engblue!10]
    plot[smooth, domain=1.45:6.0, samples=60] (\x, {0.6 + 6/(\x+0.2)})
    -- (6.0, 4.3) -- (1.45, 4.3) -- cycle;
\end{scope}

% The Pareto front: the lower-left boundary (convex, decreasing).
\draw[engred, very thick, smooth, samples=80, domain=1.5:5.9]
  plot (\x, {0.6 + 6/(\x+0.2)});
\node[engred, anchor=north west, font=\scriptsize] at (2.0, 1.5)
  {Pareto front};

% A dominated interior point.
\fill[enggray] (3.7, 3.2) circle (2pt);
\node[enggray, anchor=south west, font=\scriptsize] at (3.75, 3.25)
  {dominated};
% Arrows: it is dominated by front points (less of both objectives).
\draw[->, enggray, thin] (3.7, 3.2) -- (3.7, 2.18);
\draw[->, enggray, thin] (3.7, 3.2) -- (2.45, 3.2);

% Three scalarisation optima sitting on the front.
% Front y at x=1.7: 0.6 + 6/1.9 = 3.76
\fill[engblack] (1.7, 3.76) circle (1.6pt);
\node[engblack, anchor=west, font=\scriptsize] at (1.78, 3.86)
  {$w$ small};
% Front y at x=2.8: 0.6 + 6/3.0 = 2.6
\fill[engblack] (2.8, 2.6) circle (1.6pt);
\node[engblack, anchor=south west, font=\scriptsize] at (2.88, 2.66)
  {$w$ mid};
% Front y at x=4.8: 0.6 + 6/5.0 = 1.8
\fill[engblack] (4.8, 1.8) circle (1.6pt);
\node[engblack, anchor=south west, font=\scriptsize] at (4.88, 1.86)
  {$w$ large};

\end{tikzpicture}
\caption[The Pareto front of a bi-objective design]{The Pareto front
of a two-objective design problem trading cost against mass. The shaded
region is the achievable set: every feasible design realises some
$(\text{cost}, \text{mass})$ pair inside it. The lower-left boundary,
drawn in red, is the Pareto front, the set of designs that cannot be
improved in one objective without worsening the other. The grey
interior point is dominated: a front design exists with both lower cost
and lower mass, so no rational engineer would choose the interior
point. Scalarising the two objectives into a single weighted sum and
minimising for several weights $w$ recovers individual front points,
black markers, with small $w$ favouring low cost and large $w$
favouring low mass. The weighted-sum sweep traces the convex part of
the front; non-convex pieces of a front require other scalarisations,
which is the reason the engineer reports the whole front rather than a
single scalar optimum.}
\label{fig:vol02:ch15:pareto-front}
\end{figure}
